% PM apres MAP
% Images pour montrer staircasing



\section{Introduction} \label{sec:intro}
Image denoising problems consist in estimating an unknown image from a noisy observation in a way that accounts for the underlying uncertainties. Variational and Bayesian methods have become two important approaches for doing so, but the question of which of these two approaches is preferable to apply for a given imaging problem remains contested \cite{louchet2008modeles,louchet2013posterior,burger2014maximum,gribonval2018bayesian}. As variational methods correspond to maximum likelihood estimation in a Bayesian setting, this question is often centered around the merits and shortcomings of maximum a posteriori (MAP) estimators over Bayesian posterior mean (PM) estimators for image denoising problems. The goal of this paper is to describe a broad class of Bayesian posterior mean estimators with Gaussian data fidelity term and log-concave prior using Hamilton--Jacobi (HJ) partial differential equations (PDEs) and use these connections to clarify certain image denoising properties of this class of Bayesian posterior estimators.

To illustrate the main ideas of this paper, we first briefly introduce convex finite-dimensional variational methods relevant to image denoising problems and describe their connections to first-order HJ PDEs with initial data. We then discuss how these ideas can be extended to a broad class of Bayesian methods and posterior mean estimators.

Variational methods formulate image denoising problems as the optimization of a weighted sum of a data fidelity term (which models the noise corrupting the unknown image) and a regularization term (which models known properties of the unknown image) \cite{chambolle2016introduction,darbon2015convex}, where the goal is to minimize this sum to obtain an estimate that hopefully accounts well for both the data fidelity term and the regularization term. This approach is popular because the resultant optimization problem for various non-smooth and convex regularization terms used in image denoising problems, such as total variation and $l_1$-norm based regularization terms, is well-understood \cite{bouman1993generalized,candes2006robust,chambolle1997image,darbon2015convex,darbon2019decomposition,donoho2006compressed,rudin1992nonlinear} and can be solved efficiently using robust numerical optimization methods \cite{chambolle2016introduction}. These methods may have shortcomings, however, in that reconstructed images from variational methods with non-smooth and convex regularization terms may have undesirable and visually unpleasant staircasing effects due to the singularities of the non-smooth regularization terms \cite{nikolova2007model,louchet2008modeles,woodford2009global}.

We will focus in this paper on the following class of variational imaging models: given an observed image $\bx\in\Rn$, where $n \in \N$ corresponds to the number of pixels in the image, corrupted by additive Gaussian noise and a parameter $t>0$, estimate the original uncorrupted image by solving the minimization problem
\begin{equation} \label{eq:minimization_prob} 
S_{0}(\bx,t)\coloneqq\inf_{\by\in\Rn}\left\{ \frac{1}{2t}\normsq{\bx-\by}+J(\by)\right\}
\end{equation}
and computing its minimizer
\begin{equation} \label{eq:minimizer_prob} 
\bu_{MAP}(\bx,t) \coloneqq \argmin_{\by\in\Rn}\left\{ \frac{1}{2t}\normsq{\bx-\by}+J(\by)\right\}. 
\end{equation}
The functions $\by\mapsto \frac{1}{2t}\normsq{\bx-\by}$ and $J\colon\Rn\to\R\cup\{+\infty\}$ in \eqref{eq:minimization_prob} correspond, respectively, to the data fidelity term and regularization of this class of variational models. The parameter $t>0$ controls the relative importance of the data fidelity term over the regularization term in \eqref{eq:minimization_prob}. When $J$ is proper, lower semicontinuous and convex (see Sect.~\ref{sec:background} for the definitions), the solution $(\bx,t)\mapsto S(\bx,t)$ to the minimization problem \eqref{eq:minimization_prob} solves the first-order HJ equation with initial data (see Theorem~\ref{thm:e_u_nonvisc_hj} in Sect.~\ref{sec:background}, definition~\ref{def:moreau_yosida})
\begin{equation}
\begin{dcases}
\frac{\partial S_0}{\partial t}(\bx,t)+\frac{1}{2}\normsq{\nabla_{\bx}S_0(\bx,t)}=0 & \mbox{in }\Rn\times(0,+\infty),\\
S_0(\bx,0)=J(\bx) & \mbox{in }\Rn.
\end{dcases}
\end{equation}
The properties of the minimization problem~\eqref{eq:minimization_prob} follow from the properties of the solution to this HJ equation~\cite{darbon2015convex}; the minimizer $\bu_{MAP}(\bx,t)$, in particular, satisfies the representation formula
\begin{equation}
\bu(\bx,t)=\bx-t\nabla_{\bx}S_0(\bx,t).
\end{equation}

The main idea of this paper is to interpret the minimization problem~\eqref{eq:minimization_prob} in a Bayesian framework as the mode of an appropriate posterior distribution, use this distribution to generate the corresponding Bayesian posterior mean estimator, and finally relate the Bayesian posterior mean estimator to the solution to a viscous HJ PDE. Specifically, given a smoothing parameter $\epsilon > 0$, and under appropriate conditions on $J$, we formally interpret the minimizer~\eqref{eq:minimizer_prob} as the mode or MAP estimate of the Bayesian posterior distribution
\begin{equation*}
\bu_{MAP}(\bx,t) = \argmin_{\by\in\Rn}\left\{ \frac{1}{2t}\left\Vert \bx-\by\right\Vert _{2}^{2}+J(\by)\right\} \equiv \argmax_{\by\in\Rn}\left\{\frac{e^{-\left(\frac{1}{2t}\left\Vert \bx-\by\right\Vert _{2}^{2}+J(\by)\right)/\epsilon}}{\int_{\Rn}e^{-\left(\frac{1}{2t}\left\Vert \bx-\by\right\Vert _{2}^{2}+J(\by)\right)/\epsilon}\diff\by}\right\}.
\end{equation*}
The posterior distribution gives the Bayesian posterior mean estimate in terms of the observed image $\bx$ and parameters $t$ and $\epsilon$:
\begin{equation} \label{eq:pm_estimate}
    \bu_{PM}(\bx,t,\epsilon) \coloneqq \frac{\int_{\Rn} \by e^{-\left(\frac{1}{2t}\normsq{\bx-\by} + J(\by)\right)/\epsilon} \diff\by}{\int_{\Rn} e^{-\left(\frac{1}{2t}\normsq{\bx-\by} + J(\by)\right)/\epsilon} \diff\by}.
\end{equation}
We will show that under appropriate conditions on a generalized prior $e^-{J}$, the posterior mean estimate~\eqref{eq:pm_estimate} is described by the solution to a viscous HJ PDE with initial data $J$ similarly to how the MAP estimate~\ref{eq:minimizer_prob} is described by the solution to the first-order HJ PDE~\eqref{eq:hj_non_visc_pde}. Specifically, the function $S_\epsilon\colon\Rn\times[0,+\infty)\mapsto[0,+\infty)$ defined by
\begin{equation}\label{eq:intro_Seps}
S_{\epsilon}(\bx,t) \coloneqq -\epsilon\ln\left(\frac{1}{\left(2\pi t\epsilon\right)^{n/2}}\int_{\Rn}e^{-\left(\frac{1}{2t}\left\Vert \bx-\by\right\Vert _{2}^{2}+J(\by)\right)/\epsilon}\diff\by\right)
\end{equation}
solves the viscous HJ equation with initial data
\begin{equation}
\begin{dcases}\label{intro:pde}
\frac{\partial S_{\epsilon}}{\partial t}(\bx,t)+\frac{1}{2}\left\Vert \nabla_{\bx}S_{\epsilon}(\bx,t)\right\Vert _{2}^{2}=\frac{\epsilon}{2}\Laplacian S_{\epsilon}(\bx,t) & \text{ in }\Rn\times(0,+\infty),\\
S_{\epsilon}(\bx,0)=J(\bx) & \text{ in }\Rn,
\end{dcases}
\end{equation}
and the PM estimate satisfies the representation formula
\begin{equation}
\bu_{PM}(\bx,t,\epsilon)=\bx-t\nabla_{\bx}S_{\epsilon}(\bx,t).
\end{equation}
Various properties of PM estimators will follow from the properties to both the solutions $S_0(\bx,t)$ and $S_{\epsilon}(\bx,t)$ of these first-order and viscous HJ PDEs.

The main motivations for this paper are twofold. The first motivation comes from the papers of \cite{louchet2008modeles,gribonval2011should,louchet2013posterior}, which showed that the class of Bayesian posterior mean estimators with Gaussian data fidelity term \eqref{eq:pm_estimate} and certain priors can be expressed as minimizers to optimization problems involving a quadratic fidelity term and a smooth convex regularization term \cite{louchet2008modeles,gribonval2011should,gribonval2013reconciling,gribonval2018characterization,gribonval2018bayesian}, i.e., there exists a smooth regularization term $f_{\rm{prior}}\colon \Rn \mapsto \R$ such that
\[
\bu_{PM}(\bx,t,\epsilon) = \argmin_{\by \in \Rn} \left\{\frac{1}{2}\normsq{\bx-\by} + f_{\rm{prior}}(\by)\right\}.
\]
The second motivation comes from the papers of \cite{burger2014maximum,burger2016bregman,pereyra2019revisiting} who showed that the MAP estimate can under certain assumptions on the prior $J$ be characterized as a proper Bayes estimator, in that it minimizes in expectation a loss function.

% (and later by \cite{gribonval2018characterization,gribonval2018bayesian} to certain non-Gaussian data fidelity terms)

To our knowledge, \cite{louchet2008modeles,louchet2013posterior} derived this result when the regularization term $J$ in the posterior mean estimate~\eqref{eq:pm_estimate} is a total variation term~\cite{rudin1992nonlinear}. This result was later generalized to a broader class of regularization terms by \citet{gribonval2011should}, and to colored Gaussian and other non-Gaussian data fidelity terms in \cite{gribonval2013reconciling,gribonval2018bayesian}. All these studies, however, only proved existence of the regularization term $f_{\rm{prior}}$, however, and not their explicit form. The connections between HJ PDEs and posterior mean estimators that we derive in this paper will allow us to find the representation of $f_{\rm{prior}}$, and moreover establish that these are twice continuously differentiable and satisfy Theorem 3 of \citet{nikolova2004weakly}, thereby proving that these Bayesian posterior mean estimators avoid image denoising staircasing effects.

In addition to its relevance to image denoising problems, the viscous HJ PDE~\eqref{intro:pde} has recently received some attention in the deep learning literature, where its solution $\bx \mapsto S_\epsilon(\bx,t)$ is known as the local entropy loss function and is a loss regularization effective at training deep networks \cite{chaudhari2018deep,chaudhari2019entropy,garcia2019variational,vidal2017mathematics}. While this paper focuses on HJ PDEs and Bayesian posterior mean estimates in imaging sciences, the results in this paper should also be relevant to the deep learning literature and give new theoretical understandings of the local entropy loss function in terms of the data $\bx$ and parameters $t$ and $\epsilon$.

Now to illustrate these properties, we give an example based on the Rudin--Osher--Fatemi (ROF) image denoising model, which consists of a total variation (TV) prior with quadratic data fidelity term \cite{bouman1993generalized,chambolle1997image,rudin1992nonlinear}. The finite dimension anisotropic TV prior endowed with 4-nearest neighbors interactions \cite{winkler2003image}, for instance, takes the form
\begin{equation*}
    TV(\by) = \frac{1}{2}\sum_{i \in \mathcal{V}}\sum_{j \in \mathcal{N}(i)}|y_i-y_j|,
\end{equation*}
where the lattice $\mathcal{V}$ corresponds to a regular two-dimensional grid, and for every $i$ we denote by $\mathcal{N}(i)$ the set of the 4-nearest neighbors of $i$. Letting the observed image $\bx$ is defined on this lattice, with $|\mathcal{V}|= n$, the associated anisotropic ROF problem \cite{rudin1992nonlinear} then takes the form
\begin{equation}\label{eq:rof_model}
\inf_{y_i \in \mathbb{R} \colon i \in \mathcal{V}}\left\{\sum_{i \in \mathcal{V}} \frac{1}{2t}(x_i - y_i)^2 + \frac{1}{2}\sum_{i \in \mathcal{V}}\sum_{j \in \mathcal{N}(i)}|y_i-y_j| \right\} \equiv \inf_{\by \in \Rn}\left\{\frac{1}{2t}\normsq{\bx-\by} + TV(\by) \right\}.
\end{equation}
In a Bayesian setting, the anisotropic ROF problem above corresponds to the mode of the Bayesian posterior distribution:
\begin{equation}\label{eq:rof_dist}
\by \mapsto \frac{e^{-\left(\frac{1}{2t}\normsq{\bx-\by} + TV(\by)\right)/\epsilon}}{\int_{\Rn} e^{-\left(\frac{1}{2t}\normsq{\bx-\by} + TV(\by)\right)/\epsilon} \diff\by}
\end{equation}
and the corresponding posterior mean estimate is
\begin{equation}\label{eq:pm_tv}
\bu_{PM}(\bx,t,\epsilon) = \frac{\int_{\Rn} \by e^{-\left(\frac{1}{2t}\normsq{\bx-\by} + TV(\by)\right)/\epsilon} \diff \by}{\int_{\Rn} e^{-\left(\frac{1}{2t}\normsq{\bx-\by} + TV(\by)\right)/\epsilon} \diff\by}.
\end{equation}

\begin{figure}[h]
    \centering
\begin{subfigure}{0.55\textwidth}
    \centering
    \includegraphics[width = 0.60\textwidth]{Images/fig2_barbara_original.png}
    \caption{}
    \label{subfig:lena}
\end{subfigure}%
\begin{subfigure}{0.55\textwidth}
    \centering
    \includegraphics[width = 0.60\textwidth]{Images/fig2_barbara_noisy.png}
    \caption{}
    \label{subfig:lena_noisy}
\end{subfigure}

\begin{subfigure}{0.55\textwidth}
    \centering
    \includegraphics[width = 0.60\textwidth]{Images/fig2_barbara_u_map.png}
    \caption{}
    \label{subfig:lena_denoised_umap}
\end{subfigure}%
\begin{subfigure}{0.55\textwidth}
    \centering
    \includegraphics[width = 0.60\textwidth]{Images/fig2_barbara_u_pm_m=100000.png}
    \caption{}
    \label{subfig:lena_denoised_upm}
\end{subfigure}
    \caption{The anisotropic ROF model for image denoising applied to the noisy grayscale 256x256 image of Barbara.  The noise is Gaussian with zero mean and standard deviation $\sigma = 20$, and the parameters of the ROF model~\eqref{eq:rof_model} and Bayesian posterior distribution~\eqref{eq:rof_dist} are $t=20$ and $\epsilon=20$. (A) Original image of Barbare (B) Noisy image of Barbara (C) Denoised image of Barbara with the MAP estimate (D) Denoised image of Barbara with PM estimate.}
    \label{fig:original_images}
\end{figure}

Figure~\ref{fig:original_images}(A) depicts the standard Barbara picture (grayscale, 256x256), which we corrupt with Gaussian noise (zero mean with standard deviation $\sigma = 20$ with pixel values in $[0,255]$) in Figure~\ref{fig:original_images}(B). The noisy image is denoised by computing the minimizer $\bu_{MAP}(\bx,t,\epsilon)$ to the ROF model~\ref{eq:rof_model} and illustrated in Figure~\ref{fig:original_images}(C), and the posterior mean estimate $\bu_{PM}(\bx,t,\epsilon)$ of the Bayesian posterior distribution~\ref{eq:rof_dist} with $t = \epsilon = 20$ and illustrated in Figure~\ref{fig:original_images}. To compute the posterior mean estimate, we implemented the denoising algorithm outlined in Chapter 2 of \cite{louchet2008modeles} with convergence criterion outlined by Theorem 2.5 in Section 2.3 of \cite{louchet2008modeles}. In the terminology of their paper, the number of iterations was 100000, the subsampling rate is $256^2$, and the perturbation size $\alpha$ is approximately 50. Figure 1(C) illustrates the denoised image using the minimizer to the TV problem~\eqref{eq:rof_model}, and Figure 1(D) illustrates the denoised image using the posterior mean estimate~\eqref{eq:pm_tv}. The denoised image in Figure 1(C) contains regions where the pixels are equal and lead visually to staircasing effects. In contrast, the denoised image in Figure 1(D) does not contain such regions. This property is not an accident; it was demonstrated in \citet{louchet2008modeles} that the posterior mean estimate with TV prior does not lead to staircasing effects. In Section~\ref{subsec:HJ_eqns_1order}, we will show in fact that the posterior mean estimate~\ref{eq:pm_estimate} avoids staircasing effects for a broad class of convex priors $J$.

\subsection{Contributions} \label{subsec:intro_contributions}
In this paper, we establish connections between a broad class of Bayesian methods and posterior mean estimators and solutions to HJ PDEs. These connections are described in Theorems~\ref{thm:visc_hj_eqns} and~\ref{thm:char_opt_sol} for viscous HJ PDEs and first-order HJ PDEs, respectively. We show in Theorem~\ref{thm:visc_hj_eqns} that the posterior mean estimate $\bu_{PM}(\bx,t,\epsilon)$ is described by the solution to the viscous HJ PDE~\eqref{intro:pde} with initial data corresponding to the convex prior $J$, which we characterize in detail in terms of the data $\bx$ and parameters $t$ and $\epsilon$. In particular, the posterior mean estimate satisfies the representation formula $\bu_{PM}(\bx,t,\epsilon) = \bx - t\nabla_{\bx}S_\epsilon(\bx,t)$. Next, we use the connections between viscous HJ PDEs and posterior mean estimates established in Theorem~\ref{thm:visc_hj_eqns} to show in Theorem~\ref{thm:char_opt_sol} that the posterior mean estimate $\bu_{PM}(\bx,t,\epsilon)$ can be expressed through the gradient of the solution to a first-order HJ PDE with twice continuously differentiable convex initial data $\Rn \ni \bx \mapsto K^{*}_{\epsilon}(\bx,t)-\frac{1}{2}\Vert \bx \Vert_{2}^{2}$, where
\[
K_{\epsilon}(\bx,t) = t\epsilon\ln\left(\frac{1}{(2\pi t\epsilon)^{n/2}}\int_{\dom J}e^{\left(\frac{1}{t}\left\langle \bx,\by\right\rangle -\frac{1}{2t}\left\Vert \by\right\Vert _{2}^{2}-J(\by)\right)/\epsilon}\diff\by\right)
\]
and $\bx \mapsto  K^{*}_{\epsilon}(\bx,t)$ is the Fenchel--Legendre transform of the function $\bx \mapsto K_{\epsilon}(\bx,t)$. In other words, we show
\[
\bu_{PM}(\bx,t,\epsilon) = \argmin_{\by \in \Rn}\left\{\frac{1}{2}\normsq{\bx-\by} + \left( K^{*}_{\epsilon}(\by,t)-\frac{1}{2}\Vert \by \Vert_{2}^{2}\right)\right\}.
\]
This formula gives the representation of the convex prior enabling one to express the posterior mean estimate as the minimizer of a convex variational problem, and in fact in terms of the solution to a first-order HJ PDE, thereby extending the results of \cite{louchet2008modeles,gribonval2011should} who showed existence of this prior when the data fidelity term is Gaussian, but not its representation. The twice continuously differentiability of this prior, in particular, implies that the posterior mean estimate $\bu_{PM}(\bx,t,\epsilon)$ avoids image denoising staircasing effects as a consequence of the results derived by \citet{nikolova2004weakly} (specifically, by Theorem 3 in her paper).

We also use the connections between posterior mean estimators and solutions to viscous HJ PDEs established in Theorem~\ref{thm:visc_hj_eqns} to prove several topological, representation, and monotonicity properties of posterior mean estimators, respectively, in Propositions~\ref{prop:topo_properties}, \ref{prop:pm_rep_props}, and \ref{prop:mono_properties}. These properties are then used in Proposition\ref{prop:bounds_limits} to derive an optimal upper bound on the minimum mean square (MMS) error $\expectationJ{\normsq{\by-\bu_{PM}(\bx,t,\epsilon)}}$, several estimates on the MAP and posterior mean estimates, and the behavior of the posterior mean estimate $\bu_{PM}(\bx,t,\epsilon)$ in the limit $t \to 0$. 
Finally, we use the connections between both MAP and posterior mean estimates and HJ PDEs to characterize the MAP and posterior mean estimates in the context of Bayesian estimation theory, and specifically in Theorem~\ref{thm:bregman_div} to show that the MAP estimate~\eqref{eq:minimizer_prob} corresponds to the Bayes estimator of the Bayesian risk~\eqref{eq:breg2} whenever the prior $J$ is convex on $\Rn$. This result, in particular, extends the findings of \cite{burger2014maximum} and \cite{pereyra2019revisiting} by removing the restriction on the prior $J$ to be uniformly Lipschitz continuous on $\Rn$, at least when the data fidelity term is Gaussian. This result does not generalize to the case where $J$ is not defined everywhere on $\Rn$, but nonetheless we show that in this situation that the Bayesian risk~\eqref{eq:breg2} is a Bayes estimator described in terms of the solution to both the first-order HJ PDE~\eqref{thm:e_u_nonvisc_hj} and the viscous HJ PDE~\eqref{thm:visc_hj_eqns}.

\subsection{Organization}\label{subsec:intro_organization}
In Section \ref{sec:background}, we review concepts of real and convex analysis that will be used throughout this paper. In Section~\ref{sec:connections}, we establish theoretical connections between a broad class of Bayesian posterior mean estimators and HJ PDEs. Our mathematical set-up is described in Subsection~\ref{subsec:set-up}, the connections of posterior mean estimators to viscous HJ PDEs are described in Subsection~\ref{subsec:HJ_eqns_visc}, and the connections of posterior mean estimators to first-order HJ PDEs are described in Subsection~\ref{subsec:HJ_eqns_1order}. We use these connections to establish various properties of posterior mean estimators in Section~\ref{sec:properties_estimators}. Specifically, we establish topological, representation, and monotonicity properties of posterior mean estimators in Subsection~\ref{subsec:prop_topo-rep-mono}, an optimal upper bound for the MMS error, various estimates and bounds involving the posterior mean estimate, and the behavior of the posterior mean estimate in the limit $t \to 0$ in Subsection~\ref{subsec:prop_bounds-est}. Finally, we establish properties of MAP and posterior mean estimators in terms of Bayesian risks involving Bregman divergences in Subsection~\ref{subsec:prop_breg_div}.